\documentclass{ximera}

\begin{document}

    \title{Lecture 8}
    
    \begin{abstract}  
    Outline: \\
    - Cauchy sequences
\end{abstract}

\maketitle

\begin{definition}
    A sequence $\{a_n\}$ is called a Cauchy sequence if for every $\epsilon > 0$, there exists an $n \in \mathbb{N}$ such that for any $m,n\geq N$, $|a_n-a_m|< \epsilon$.
\end{definition}

\begin{exercise}
    Prove that every convergent sequence is Cauchy.
\end{exercise}

\begin{exercise}
    Come up with an example of a Cauchy sequence in the rationals that is not convergent in the rationals.
\end{exercise}

\begin{exercise}
    Prove Cauchy sequences are bounded.
\end{exercise}
\begin{proof}
    Taking $\epsilon = 1$, there exists $N$ such that $|x_n - x_N| < 1$ for all $n \geq N$.  Then reverse triangle inequality says $||x_n| - |x_N|| \leq |x_n - x_N| < 1$, and therefore $|x_n| < |x_N| + 1$.  Thus, an upper bound is
    \begin{align}
        M = \max\{|x_1|, ..., |x_N|, |x_N|+1\}.
    \end{align}
\end{proof}

\end{document}